%%%
%%% Radical "⽕"
%%%
\section*{Radical 86: ``⽕'' (灬)}\addcontentsline{toc}{section}{Radical 86: ⽕,灬}\addcontentsline{loh}{figure}{\#\#\#\# 86: ⽕}

%%%%%%%%%% 火 %%%%%%%%%%
\subsection*{火}\addcontentsline{loh}{figure}{火}

\begin{Entry}{火}{4}{⽕}[Kangxi 86]
  \begin{Phonetics}{火}{huo3}[][HSK 3,4]
    \definition*{s.}{Sobrenome: Huo}
    \definition{adj.}{ardente; flamejante; vermelho como fogo | efervescente; próspero}
    \definition{adv.}{urgentemente}
    \definition{clas.}{usado para unidades militares (antigo)}
    \definition[场,把,团,堆]{s.}{fogo; a luz e as chamas emitidas pela combustão de um objeto | fúria; metáfora para emoções agitadas, irritadas ou raivosas | calor interno (uma das seis causas de doenças) | armas de fogo e munições | a ação de lutar}
    \definition{v.}{ficar com raiva; perder a paciência}
  \end{Phonetics}
\end{Entry}

\begin{Entry}{火山}{4,3}{⽕,⼭}
  \begin{Phonetics}{火山}{huo3shan1}[][HSK 7-9]
    \definition[座]{s.}{vulcão}
  \end{Phonetics}
\end{Entry}

\begin{Entry}{火车}{4,4}{⽕,⾞}
  \begin{Phonetics}{火车}{huo3che1}[][HSK 1]
    \definition[个,列,节,班,趟]{s.}{trem; comboio}
  \antonymref{列车}{lie4che1}
  \end{Phonetics}
\end{Entry}

\begin{Entry}{火车司机}{4,4,5,6}{⽕,⾞,⼝,⽊}
  \begin{Phonetics}{火车司机}{huo3che1 si1ji1}
    \definition{s.}{maquinista; engenheiro (de locomotiva)}
  \end{Phonetics}
\end{Entry}

\begin{Entry}{火灾}{4,7}{⽕,⽕}
  \begin{Phonetics}{火灾}{huo3zai1}[][HSK 5]
    \definition[起,场]{s.}{fogo (como um desastre); conflagração; desastres causados por incêndios}
  \end{Phonetics}
\end{Entry}

\begin{Entry}{火花}{4,7}{⽕,⾋}
  \begin{Phonetics}{火花}{huo3hua1}[][HSK 7-9]
    \definition[簇]{s.}{faísca; explosão de chamas | um padrão brilhante}
  \end{Phonetics}
\end{Entry}

\begin{Entry}{火炬}{4,8}{⽕,⽕}
  \begin{Phonetics}{火炬}{huo3ju4}[][HSK 7-9]
    \definition[把]{s.}{tocha}
  \end{Phonetics}
\end{Entry}

\begin{Entry}{火药}{4,9}{⽕,⾋}
  \begin{Phonetics}{火药}{huo3yao4}[][HSK 7-9]
    \definition[桶,克]{s.}{pólvora; um tipo de explosivo que explode com fumaça, como pólvora preta, ou sem fumaça, como nitrato de celulose}
  \end{Phonetics}
\end{Entry}

\begin{Entry}{火候}{4,10}{⽕,⼈}
  \begin{Phonetics}{火候}{huo3hou5}[][HSK 7-9]
    \definition{s.}{duração e grau de aquecimento, cozimento, fundição, etc.  |Coloquial: nível de realização | Coloquial: momento crucial, crítico; emergência}
  \end{Phonetics}
\end{Entry}

\begin{Entry}{火柴}{4,10}{⽕,⽊}
  \begin{Phonetics}{火柴}{huo3chai2}[][HSK 5]
    \definition[根,盒,包]{s.}{fósforo (palito de fósforo); fósforo de segurança; iniciador de fogo feito de uma tira fina de madeira mergulhada em um composto de fósforo ou enxofre}
  \end{Phonetics}
\end{Entry}

\begin{Entry}{火海}{4,10}{⽕,⽔}
  \begin{Phonetics}{火海}{huo3hai3}
    \definition{s.}{mar de chamas}
  \antonymref{雪山}{xue3shan1}
  \end{Phonetics}
\end{Entry}

\begin{Entry}{火热}{4,10}{⽕,⽕}
  \begin{Phonetics}{火热}{huo3re4}[][HSK 7-9]
    \definition{adj.}{ardente; fervente; fervoroso; apaixonado; escaldante; abrasador (opp. 冰冷)}
  \seealsoref{冰冷}{bing1leng3}
  \end{Phonetics}
\end{Entry}

\begin{Entry}{火速}{4,10}{⽕,⾡}
  \begin{Phonetics}{火速}{huo3su4}[][HSK 7-9]
    \definition{adv.}{em alta velocidade; com pressa; usar a velocidade mais rápida (para fazer coisas urgentes)}
  \end{Phonetics}
\end{Entry}

\begin{Entry}{火焰}{4,12}{⽕,⽕}
  \begin{Phonetics}{火焰}{huo3yan4}[][HSK 7-9]
    \definition[团,缕,股,道]{s.}{chama; labareda; flama}
  \end{Phonetics}
\end{Entry}

\begin{Entry}{火锅}{4,12}{⽕,⾦}
  \begin{Phonetics}{火锅}{huo3guo1}[][HSK 7-9]
    \definition[顿]{s.}{\emph{hot pot}; uma panela feita de metal ou outro material, que pode ser usada para aquecer a sopa continuamente com eletricidade, álcool, etc., e depois adicionar carne, vegetais, etc. à sopa e comê"-la enquanto cozinha; atualmente, refere"-se principalmente a alimentos cozidos dessa maneira}
  \end{Phonetics}
\end{Entry}

\begin{Entry}{火腿}{4,13}{⽕,⾁}
  \begin{Phonetics}{火腿}{huo3tui3}[][HSK 5]
    \definition[道,个]{s.}{presunto; as pernas de porco marinadas mais famosas são produzidas em Jinhua, na província de Zhejiang, e em Xuanwei, na província de Yunnan.}
  \end{Phonetics}
\end{Entry}

\begin{Entry}{火辣辣}{4,14,14}{⽕,⾟,⾟}
  \begin{Phonetics}{火辣辣}{huo3la4la4}[][HSK 7-9]
    \definition{adj.}{ardente; escaldante; abrasador}
  \end{Phonetics}
\end{Entry}

\begin{Entry}{火暴}{4,15}{⽕,⽇}
  \begin{Phonetics}{火暴}{huo3bao4}[][HSK 7-9]
    \definition{adj.}{impetuoso; impetuoso; irritável; impaciente | próspero; florescente; vigoroso; animado}
  \end{Phonetics}
\end{Entry}

\begin{Entry}{火箭}{4,15}{⽕,⽵}
  \begin{Phonetics}{火箭}{huo3jian4}[][HSK 6]
    \definition[个,艘,发,枚]{s.}{foguete; uma aeronave de alta velocidade que utiliza força de reação para se impulsionar para a frente; é usado para lançar satélites, naves espaciais, etc.; também pode ser equipado com uma ogiva para fabricar um míssil}
  \end{Phonetics}
\end{Entry}

%%%%%%%%%% 灭 %%%%%%%%%%
\subsection*{灭}\addcontentsline{loh}{figure}{灭}

\begin{Entry}{灭}{5}{⽕}
  \begin{Phonetics}{灭}{mie4}[][HSK 6]
    \definition{v.}{extinguir"-se | extinguir; apagar; desligar | afogar; inundar; submergir | perecer; destruir | exterminar; apagar; acabar com; tornar inexistente}
  \end{Phonetics}
\end{Entry}

\begin{Entry}{灭亡}{5,3}{⽕,⼇}
  \begin{Phonetics}{灭亡}{mie4wang2}[][HSK 7-9]
    \definition{v.}{ser destruído; extinguir"-se; perecer; desaparecer; isso se refere à eliminação de uma nação, grupo étnico ou grupo político, que deixa de existir | destruir; exterminar}
  \end{Phonetics}
\end{Entry}

\begin{Entry}{灭火}{5,4}{⽕,⽕}
  \begin{Phonetics}{灭火}{mie4huo3}
    \definition{s.}{combate a incêndios}
    \definition{v.}{apagar um incêndio; extinguir um incêndio | desligar o motor (de combustão)}
  \antonymref{点火}{dian3/huo3}
  \antonymref{燃烧}{ran2shao1}
  \end{Phonetics}
\end{Entry}

\begin{Entry}{灭绝}{5,9}{⽕,⽷}
  \begin{Phonetics}{灭绝}{mie4jue2}[][HSK 7-9]
    \definition{v.}{extinguir"-se; ser extinto; eliminar completamente | perder completamente; perder totalmente}
  \end{Phonetics}
\end{Entry}

%%%%%%%%%% 灯 %%%%%%%%%%
\subsection*{灯}\addcontentsline{loh}{figure}{灯}

\begin{Entry}{灯}{6}{⽕}
  \begin{Phonetics}{灯}{deng1}[][HSK 2]
    \definition*{s.}{Sobrenome: Deng}
    \definition[盏,个]{s.}{lâmpada; luz; lanterna; dispositivo luminoso, usado principalmente para iluminação | queimador; um aparelho que brilha e aquece como uma lâmpada e pode ser usado para aquecer | tubo; válvula; o nome popular dado aos tubos eletrônicos com formato semelhante a lâmpadas encontrados em aparelhos antigos, como rádios}
  \seealsoref{灯火}{deng1huo3}
  \synonymref{灯火}{deng1huo3}
  \end{Phonetics}
\end{Entry}

\begin{Entry}{灯火}{6,4}{⽕,⽕}
  \begin{Phonetics}{灯火}{deng1huo3}
    \definition{s.}{luzes; iluminação}
  \seealsoref{灯}{deng1}
  \synonymref{灯}{deng1}
  \synonymref{灯光}{deng1guang1}
  \end{Phonetics}
\end{Entry}

\begin{Entry}{灯丝}{6,5}{⽕,⼀}
  \begin{Phonetics}{灯丝}{deng1si1}
    \definition[根]{s.}{filamento (em uma lâmpada ou válvula); ligamentos; o filamento metálico dentro de uma lâmpada ou tubo eletrônico é geralmente um filamento fino de tungstênio, que pode emitir luz, calor, elétrons ou raios quando uma corrente elétrica passa por ele}
  \end{Phonetics}
\end{Entry}

\begin{Entry}{灯号}{6,5}{⽕,⼝}
  \begin{Phonetics}{灯号}{deng1hao4}
    \definition{s.}{sinal luminoso | luz intermitente | luz indicadora}
  \synonymref{信号}{xin4hao4}
  \end{Phonetics}
\end{Entry}

\begin{Entry}{灯光}{6,6}{⽕,⼉}
  \begin{Phonetics}{灯光}{deng1guang1}[][HSK 4]
    \definition[束,盏,点,打]{s.}{iluminação; luminosidade da lâmpada | luminação (palco); equipamento de iluminação para palco ou estúdio}
  \end{Phonetics}
\end{Entry}

\begin{Entry}{灯泡}{6,8}{⽕,⽔}
  \begin{Phonetics}{灯泡}{deng1pao4}[][HSK 7-9]
    \definition[只,个]{s.}{lâmpada (bulbo) | (gíria) terceiro indesejado estragando encontro de casal; é frequentemente usado para descrever a si mesmo ou a outros se sentindo estranhos ou indesejados em situações sociais}
  \seealsoref{电灯泡}{dian4deng1pao4}
  \end{Phonetics}
\end{Entry}

\begin{Entry}{灯标}{6,9}{⽕,⽊}
  \begin{Phonetics}{灯标}{deng1biao1}
    \definition{s.}{luz de farol; farol; bóia luminosa}
  \end{Phonetics}
\end{Entry}

\begin{Entry}{灯笼}{6,11}{⽕,⽵}
  \begin{Phonetics}{灯笼}{deng1long5}[][HSK 7-9]
    \definition[个,盏,只]{s.}{lanterna; luminárias suspensas ou portáteis, geralmente feitas de finas tiras de bambu ou arame de ferro como estrutura, cobertas com areia ou papel, com velas dentro; atualmente, lâmpadas elétricas são usadas principalmente como fontes de luz e como decoração}
  \end{Phonetics}
\end{Entry}

%%%%%%%%%% 灰 %%%%%%%%%%
\subsection*{灰}\addcontentsline{loh}{figure}{灰}

\begin{Entry}{灰}{6}{⽕}
  \begin{Phonetics}{灰}{hui1}[][HSK 7-9]
    \definition{adj.}{cinza (cor) | desanimado; desencorajado; deprimido}
    \definition[把,堆]{s.}{cinzas; pó que sobra após a queima de um objeto | pó; poeira; substância em pó | cal; argamassa (de cal)}
  \end{Phonetics}
\end{Entry}

\begin{Entry}{灰心}{6,4}{⽕,⼼}
  \begin{Phonetics}{灰心}{hui1/xin1}[][HSK 7-9]
    \definition{v.+compl.}{desanimar; ficar desanimado; ficar desapontado; (devido a dificuldades, fracassos) ficar deprimido}
  \end{Phonetics}
\end{Entry}

\begin{Entry}{灰尘}{6,6}{⽕,⼩}
  \begin{Phonetics}{灰尘}{hui1chen2}[][HSK 7-9]
    \definition[层,堆]{s.}{cinza; poeira; sujeira; pó}
  \end{Phonetics}
\end{Entry}

\begin{Entry}{灰色}{6,6}{⽕,⾊}
  \begin{Phonetics}{灰色}{hui1se4}[][HSK 5]
    \definition{adj.}{obscuro; ambíguo | sombrio; pessimista}
    \definition[种]{s.}{cor cinza; acinzentado}
  \end{Phonetics}
\end{Entry}

%%%%%%%%%% 灵 %%%%%%%%%%
\subsection*{灵}\addcontentsline{loh}{figure}{灵}

\begin{Entry}{灵}{7}{⽕}
  \begin{Phonetics}{灵}{ling2}[][HSK 7-9]
    \definition*{s.}{Sobrenome: Ling}
    \definition{adj.}{rápido; inteligente; afiado | eficaz; efetivo | flexível; hábil}
    \definition{s.}{espírito; alma | inteligência; mente | fada; duende; elfo | restos mortais do falecido; esquife | carro funerário; caixão ou algo relacionado aos mortos}
  \end{Phonetics}
\end{Entry}

\begin{Entry}{灵巧}{7,5}{⽕,⼯}
  \begin{Phonetics}{灵巧}{ling2qiao3}[][HSK 7-9]
    \definition{adj.}{destro; ágil; habilidoso; engenhoso; flexível e engenhoso}
  \end{Phonetics}
\end{Entry}

\begin{Entry}{灵机一动}{7,6,1,6}{⽕,⽊,⼀,⼒}
  \begin{Phonetics}{灵机一动}{ling2ji1-yi1dong4}[][HSK 7-9]
    \definition{expr.}{``Uma inspiração repentina.''; ter uma ideia genial de repente; uma ideia brilhante surge; um lampejo de inspiração; um vislumbre de genialidade; um plano surgiu para alguém; ter uma ideia repentina; ter uma inspiração repentina; ter uma ideia brilhante; ter uma ideia; com uma inspiração repentina; tomado por um impulso repentino (inspiração); ele mudou de ideia rapidamente (geralmente se referindo a encontrar uma solução na hora)}
  \end{Phonetics}
\end{Entry}

\begin{Entry}{灵活}{7,9}{⽕,⽔}
  \begin{Phonetics}{灵活}{ling2huo2}[][HSK 6]
    \definition[种,点,些]{adj.}{ágil; rápido; ligeiro; descreve a capacidade de fazer rapidamente mudanças apropriadas com base na situação ao lidar com as coisas | flexível; elástico; descreve reações rápidas, como movimentos e funções cerebrais}
  \end{Phonetics}
\end{Entry}

\begin{Entry}{灵通}{7,10}{⽕,⾡}
  \begin{Phonetics}{灵通}{ling2tong1}[][HSK 7-9]
    \definition{adj.}{Dialeto: útil; eficaz | Dialeto: ágil; flexível}
    \definition{v.}{ter acesso rápido à informação; ser bem informado | ser útil; cumprir o propósito}
  \end{Phonetics}
\end{Entry}

\begin{Entry}{灵敏}{7,11}{⽕,⽁}
  \begin{Phonetics}{灵敏}{ling2min3}[][HSK 7-9]
    \definition{adj.}{agudo; perspicaz; sensível}
  \end{Phonetics}
\end{Entry}

\begin{Entry}{灵感}{7,13}{⽕,⼼}
  \begin{Phonetics}{灵感}{ling2gan3}[][HSK 7-9]
    \definition[种,点,个]{s.}{inspiração; uma ideia criativa repentina; explosão de criatividade em empreendimento científico ou artístico}
  \end{Phonetics}
\end{Entry}

\begin{Entry}{灵魂}{7,13}{⽕,⿁}
  \begin{Phonetics}{灵魂}{ling2hun2}[][HSK 7-9]
    \definition{s.}{alma; espírito; pessoas supersticiosas acreditam que a alma é uma entidade não material ligada ao corpo humano como mestra, e que uma pessoa morre quando a alma deixa o corpo | alma; pensamento; o mundo interior, incluindo o espírito e o pensamento | alma; consciência; a compreensão e a atitude corretas de uma pessoa em relação ao que é certo e errado, bom e mau | alma; fator decisivo; metaforicamente, refere"-se à pessoa ou fator que desempenha um papel decisivo}
  \end{Phonetics}
\end{Entry}

%%%%%%%%%% 灶 %%%%%%%%%%
\subsection*{灶}\addcontentsline{loh}{figure}{灶}

\begin{Entry}{灶}{7}{⽕}
  \begin{Phonetics}{灶}{zao4}
    \definition[口,个]{s.}{fogão de cozinha; fogão de cozinha | cozinha; bagunça; cantina}
  \end{Phonetics}
\end{Entry}

\begin{Entry}{灶台}{7,5}{⽕,⼝}
  \begin{Phonetics}{灶台}{zao4tai2}
    \definition{s.}{fogão}
  \end{Phonetics}
\end{Entry}

%%%%%%%%%% 灼 %%%%%%%%%%
\subsection*{灼}\addcontentsline{loh}{figure}{灼}

\begin{Entry}{灼}{7}{⽕}
  \begin{Phonetics}{灼}{zhuo2}
    \definition{adj.}{brilhante; luminoso}
    \definition{v.}{queimar; chamuscar | grelhar; queimar; escaldar}
  \end{Phonetics}
\end{Entry}

\begin{Entry}{灼热}{7,10}{⽕,⽕}
  \begin{Phonetics}{灼热}{zhuo2re4}[][HSK 7-9]
    \definition{adj.}{extremamente quente; escaldante | extremamente quente; incandescente | abrasador; ardente; tórrido | preocupado}
  \synonymref{闷热}{men1re4}
  \end{Phonetics}
\end{Entry}

%%%%%%%%%% 灾 %%%%%%%%%%
\subsection*{灾}\addcontentsline{loh}{figure}{灾}

\begin{Entry}{灾}{7}{⽕}
  \begin{Phonetics}{灾}{zai1}[][HSK 5]
    \definition[个,场]{s.}{calamidade; desastre | infortúnio pessoal; adversidade | azar}
  \end{Phonetics}
\end{Entry}

\begin{Entry}{灾区}{7,4}{⽕,⼖}
  \begin{Phonetics}{灾区}{zai1qu1}[][HSK 5]
    \definition{s.}{área de desastre; área afetada por catástrofes}
  \end{Phonetics}
\end{Entry}

\begin{Entry}{灾害}{7,10}{⽕,⼧}
  \begin{Phonetics}{灾害}{zai1hai4}[][HSK 5]
    \definition[场,次,个]{s.}{desastre; calamidade; danos causados pela seca, inundações, pragas, granizo, guerras, etc.}
  \end{Phonetics}
\end{Entry}

\begin{Entry}{灾难}{7,10}{⽕,⾫}
  \begin{Phonetics}{灾难}{zai1nan4}[][HSK 5]
    \definition[场,次,个,种]{s.}{desastre; sofrimento; calamidade; catástrofe; danos e sofrimentos causados por desastres naturais ou guerras}
  \end{Phonetics}
\end{Entry}

%%%%%%%%%% 灿 %%%%%%%%%%
\subsection*{灿}\addcontentsline{loh}{figure}{灿}

\begin{Entry}{灿}{7}{⽕}
  \begin{Phonetics}{灿}{can4}
    \definition{adj.}{brilhante; luminoso; resplandecente; flamejante; deslumbrante}
  \end{Phonetics}
\end{Entry}

\begin{Entry}{灿烂}{7,9}{⽕,⽕}
  \begin{Phonetics}{灿烂}{can4lan4}[][HSK 7-9]
    \definition{adj.}{brilhante; glorioso; esplêndido; resplandecente; magnífico; deslumbrante}
  \end{Phonetics}
\end{Entry}

%%%%%%%%%% 炉 %%%%%%%%%%
\subsection*{炉}\addcontentsline{loh}{figure}{炉}

\begin{Entry}{炉}{8}{⽕}
  \begin{Phonetics}{炉}{lu2}
    \definition{s.}{fogão; forno; fornalha | Metalurgia: fornalha}
  \end{Phonetics}
\end{Entry}

\begin{Entry}{炉子}{8,3}{⽕,⼦}
  \begin{Phonetics}{炉子}{lu2zi5}[][HSK 7-9]
    \definition[个,只]{s.}{fogão; forno; fornalha; aparelhos ou equipamentos usados para cozinhar, ferver água, aquecer, fundir, etc.}
  \end{Phonetics}
\end{Entry}

\begin{Entry}{炉灶}{8,7}{⽕,⽕}
  \begin{Phonetics}{炉灶}{lu2zao4}[][HSK 7-9]
    \definition[个]{s.}{fogão; lareira; fogão para cozinhar; termo coletivo para fogão e lareira}
  \end{Phonetics}
\end{Entry}

%%%%%%%%%% 炎 %%%%%%%%%%
\subsection*{炎}\addcontentsline{loh}{figure}{炎}

\begin{Entry}{炎}{8}{⽕}
  \begin{Phonetics}{炎}{yan2}
    \definition{adj.}{escaldante; ardente}
    \definition{s.}{inflamação | poder; influência}
  \end{Phonetics}
\end{Entry}

\begin{Entry}{炎热}{8,10}{⽕,⽕}
  \begin{Phonetics}{炎热}{yan2re4}[][HSK 7-9]
    \definition{adj.}{ardente; escaldante; extremamente quente}
  \synonymref{火热}{huo3re4}
  \synonymref{闷热}{men1re4}
  \antonymref{冰冷}{bing1leng3}
  \antonymref{寒冷}{han2leng3}
  \antonymref{凉爽}{liang2shuang3}
  \antonymref{清凉}{qing1liang2}
  \end{Phonetics}
\end{Entry}

\begin{Entry}{炎症}{8,10}{⽕,⽧}
  \begin{Phonetics}{炎症}{yan2zheng4}[][HSK 7-9]
    \definition{s.}{inflamação}
  \end{Phonetics}
\end{Entry}

%%%%%%%%%% 炒 %%%%%%%%%%
\subsection*{炒}\addcontentsline{loh}{figure}{炒}

\begin{Entry}{炒}{8}{⽕}
  \begin{Phonetics}{炒}{chao3}[][HSK 6]
    \definition{v.}{saltear; refogar; aquecer os alimentos em uma panela e mexer repetidamente para cozinhá"-los ou secá"-los | especular (na bolsa de valores, etc.) | exagerar; dar publicidade exagerada; a fim de ampliar a influência, por meio de publicidade repetida e exagerada na mídia | demitir; despedir}
  \end{Phonetics}
\end{Entry}

\begin{Entry}{炒作}{8,7}{⽕,⼈}
  \begin{Phonetics}{炒作}{chao3zuo4}[][HSK 6]
    \definition{v.}{promover (na mídia); exagerar artificialmente e promover ou desvalorizar de forma inadequada | especular; comprar e vender frequentemente no mercado de negociação para obter lucros}
  \end{Phonetics}
\end{Entry}

\begin{Entry}{炒股}{8,8}{⽕,⾁}
  \begin{Phonetics}{炒股}{chao3/gu3}[][HSK 6]
    \definition{v.+compl.}{especular em ações; comprar e vender ações; jogar no mercado}
  \end{Phonetics}
\end{Entry}

%%%%%%%%%% 炖 %%%%%%%%%%
\subsection*{炖}\addcontentsline{loh}{figure}{炖}

\begin{Entry}{炖}{8}{⽕}
  \begin{Phonetics}{炖}{dun4}[][HSK 7-9]
    \definition{v.}{cozinhar | aquecer algo colocando o recipiente em água quente}
  \end{Phonetics}
\end{Entry}

%%%%%%%%%% 炫 %%%%%%%%%%
\subsection*{炫}\addcontentsline{loh}{figure}{炫}

\begin{Entry}{炫}{9}{⽕}
  \begin{Phonetics}{炫}{xuan4}
    \definition{adj.}{deslumbrante; ofuscante}
    \definition{v.}{exibir; ostentar}
  \end{Phonetics}
\end{Entry}

\begin{Entry}{炫耀}{9,20}{⽕,⽻}
  \begin{Phonetics}{炫耀}{xuan4yao4}[][HSK 7-9]
    \definition{v.}{ostentar; exibir; fazer uma demonstração de; fazer as pessoas perceberem que elas têm muito dinheiro, grandes habilidades ou um status elevado | deslumbrar}
  \synonymref{高调}{gao1diao4}
  \synonymref{夸耀}{kua1yao4}
  \synonymref{卖弄}{mai4nong5}
  \antonymref{谦虚}{qian1xu1}
  \antonymref{谦逊}{qian1xun4}
  \end{Phonetics}
\end{Entry}

%%%%%%%%%% 炭 %%%%%%%%%%
\subsection*{炭}\addcontentsline{loh}{figure}{炭}

\begin{Entry}{炭}{9}{⽕}
  \begin{Phonetics}{炭}{tan4}[][HSK 7-9]
    \definition*{s.}{Sobrenome: Tan}
    \definition[块,吨]{s.}{carvão; algo semelhante a carvão}
  \antonymref{冰}{bing1}
  \end{Phonetics}
\end{Entry}

%%%%%%%%%% 炮 %%%%%%%%%%
\subsection*{炮}\addcontentsline{loh}{figure}{炮}

\begin{Entry}{炮}{9}{⽕}
  \begin{Phonetics}{炮}{bao1}
    \definition{v.}{processar; o método de preparação da medicina chinesa é colocar as ervas cruas em uma panela de ferro em alta temperatura e fritá"-las até que fiquem marrons e estourem | secar alimentos pelo calor; refogar}
  \end{Phonetics}
  \begin{Phonetics}{炮}{pao2}
    \definition{v.}{Medicina tradicional chinesa: preparar a medicina chinesa assando"-a em uma panela de ferro quente até dourar e estalar}
  \end{Phonetics}
  \begin{Phonetics}{炮}{pao4}[][HSK 6]
    \definition{s.}{arma grande; canhão; peça de artilharia | fogo de artifício | buraco de explosão cheio de dinamite | canhão, uma das peças do xadrez chinês}
  \end{Phonetics}
\end{Entry}

\begin{Entry}{炮车}{9,4}{⽕,⾞}
  \begin{Phonetics}{炮车}{pao4che1}
    \definition{s.}{veículo de artilharia; tanque de guerra}
  \end{Phonetics}
\end{Entry}

%%%%%%%%%% 炸 %%%%%%%%%%
\subsection*{炸}\addcontentsline{loh}{figure}{炸}

\begin{Entry}{炸}{9}{⽕}
  \begin{Phonetics}{炸}{zha2}[][HSK 7-9]
    \definition{v.}{explodir; estourar; romper | dinamitar; bombardear; explodir; detonar com explosivos | encolerizar"-se; explodir em fúria | correr; fugir em pânico}
  \end{Phonetics}
  \begin{Phonetics}{炸}{zha4}[][HSK 6]
    \definition{v.}{fritar em gordura ou óleo | escaldar (como forma de cozinhar)}
  \end{Phonetics}
\end{Entry}

\begin{Entry}{炸药}{9,9}{⽕,⾋}
  \begin{Phonetics}{炸药}{zha4yao4}[][HSK 6]
    \definition[包,种]{s.}{explosivo; cargas explosivas; dinamite; substâncias que explodem quando aquecidas ou impactadas, produzindo grandes quantidades de energia e gases de alta temperatura, como dinamite e pólvora negra}
  \end{Phonetics}
\end{Entry}

\begin{Entry}{炸弹}{9,11}{⽕,⼸}
  \begin{Phonetics}{炸弹}{zha4dan4}[][HSK 6]
    \definition{s.}{bomba; uma arma com invólucro de ferro e explosivos dentro que explodem quando um fusível é acionado, geralmente lançada de um avião}
  \end{Phonetics}
\end{Entry}

%%%%%%%%%% 点 %%%%%%%%%%
\subsection*{点}\addcontentsline{loh}{figure}{点}

\begin{Entry}{点}{9}{⽕}
  \begin{Phonetics}{点}{dian3}[][HSK 1]
    \definition{clas.}{hora cheia | ponto, uma unidade de medida para tipos; antigamente, a contagem do tempo durante a noite era feita por turnos, sendo cada turno dividido em cinco pontos | quantidade ínfima; um pouco; um pouquinho; alguma coisa; indica uma pequena quantidade | usado para itens}
    \definition{s.}{gota (de líquido); (ponto) pequena gota de líquido | mancha; ponto; salpico; (um pouco) Um pequeno vestígio | (ponto) traço de um caractere chinês, cuja forma é ``、''  | Matemática: ponto; refere"-se a uma figura geométrica que não tem comprimento, largura ou altura, mas apenas uma posição | gongo, instrumento musical de metal | ponto decimal; refere"-se ao ponto decimal, símbolo matemático que representa os números decimais | lugar específico | lanche leve; petisco | lugar; grau; sinalização de um determinado local ou grau | hora marcada; hora regulamentar | aspecto; característica; partes ou aspectos específicos de algo | ritmo; batida}
    \definition{v.}{andar na ponta dos pés | dar uma dica, sugestão | tocar levemente com o dedo, pincel ou vara; tocar muito brevemente; passar rapidamente | acenar; baixar ligeiramente a cabeça e levantar rapidamente | gotejar; fazer cair líquido | semear em buracos; plantar com um plantador | verificar um por um | colocar um ponto; usar caneta e outras ferramentas para adicionar ideias | sugerir; indicar; dar uma dica | decorar; realçar | selecionar; escolher; especificar o que é exigido | acender; queimar; inflamar | (pedido) comer uma pequena quantidade de comida para saciar a fome}
  \seealsoref{点燃}{dian3ran2}
  \synonymref{部}{bu4}
  \synonymref{滴}{di1}
  \synonymref{点燃}{dian3ran2}
  \synonymref{方}{fang1}
  \synonymref{绘}{hui4}
  \synonymref{面}{mian4}
  \synonymref{描}{miao2}
  \synonymref{条}{tiao2}
  \synonymref{涂}{tu2}
  \synonymref{项}{xiang4}
  \antonymref{灭}{mie4}
  \antonymref{熄}{xi1}
  \end{Phonetics}
\end{Entry}

\begin{Entry}{点子}{9,3}{⽕,⼦}
  \begin{Phonetics}{点子}{dian3zi5}[][HSK 7-9]
    \definition[滴,个]{s.}{gota (de líquido); pequenas gotas | mancha; ponto; pinta | batidas de instrumentos de percussão | ponto"-chave; aspecto vital | ideia; indicador; método}
  \end{Phonetics}
\end{Entry}

\begin{Entry}{点心}{9,4}{⽕,⼼}
  \begin{Phonetics}{点心}{dian3xin1}
    \definition[块,份,袋,包,盒]{s.}{lanche; refeição leve}
  \end{Phonetics}
  \begin{Phonetics}{点心}{dian3xin5}[][HSK 7-9]
    \definition[块,份,袋,包,盒]{s.}{doces; docinhos; sobremesa; lanches leves; pequenos alimentos além de arroz e vegetais, como bolos, biscoitos, etc.}
  \end{Phonetics}
\end{Entry}

\begin{Entry}{点火}{9,4}{⽕,⽕}
  \begin{Phonetics}{点火}{dian3/huo3}[][HSK 7-9]
    \definition{s.}{ignição}
    \definition{v.+compl.}{acender; acender o fogo; acender uma luz; disparar; iniciar a combustão; acender as chamas | inflamar | causar problemas}
  \end{Phonetics}
\end{Entry}

\begin{Entry}{点击率}{9,5,11}{⽕,⼐,⽞}
  \begin{Phonetics}{点击率}{dian3ji1lv4}[][HSK 7-9]
    \definition{s.}{Internet: taxa de cliques (CTR, \emph{click"-through rate})}
  \end{Phonetics}
\end{Entry}

\begin{Entry}{点头}{9,5}{⽕,⼤}
  \begin{Phonetics}{点头}{dian3tou2}[][HSK 2]
    \definition{v.}{acenar com a cabeça; balançar a cabeça; mover ligeiramente a cabeça para baixo; indicar permissão, aprovação, compreensão ou saudação}
  \synonymref{肯定}{ken3ding4}
  \synonymref{拍板}{pai1/ban3}
  \synonymref{认可}{ren4ke3}
  \synonymref{认同}{ren4tong2}
  \antonymref{否定}{fou3ding4}
  \antonymref{拒绝}{ju4jue2}
  \antonymref{摇头}{yao2/tou2}
  \end{Phonetics}
\end{Entry}

\begin{Entry}{点名}{9,6}{⽕,⼝}
  \begin{Phonetics}{点名}{dian3ming2}[][HSK 4]
    \definition{v.}{fazer a lista de chamada; manter o controle da presença de alguém; chamar nomes para controle de presença | mencionar alguém pelo nome}
  \end{Phonetics}
\end{Entry}

\begin{Entry}{点评}{9,7}{⽕,⾔}
  \begin{Phonetics}{点评}{dian3ping2}[][HSK 7-9]
    \definition{s.}{comentário; crítica | um comentário ponto por ponto}
    \definition{v.}{comentar; fazer comentários}
  \end{Phonetics}
\end{Entry}

\begin{Entry}{点钟}{9,9}{⽕,⾦}
  \begin{Phonetics}{点钟}{dian3zhong1}
    \definition{clas.}{(indicando a hora do dia) horas; \emph{o'clock}; tempo indicado por um relógio; de acordo com um relógio; baseado em um relógio}
  \seealsoref{小时}{xiao3shi2}
  \synonymref{小时}{xiao3shi2}
  \end{Phonetics}
\end{Entry}

\begin{Entry}{点缀}{9,11}{⽕,⽷}
  \begin{Phonetics}{点缀}{dian3zhui4}[][HSK 7-9]
    \definition{v.}{adornar; embelezar; ornamentar; realçar, decorar, fazer com que pareça melhor | usar algo apenas para mostrar; decorar a fachada; adaptar"-se à ocasião}
  \end{Phonetics}
\end{Entry}

\begin{Entry}{点燃}{9,16}{⽕,⽕}
  \begin{Phonetics}{点燃}{dian3ran2}[][HSK 5]
    \definition{v.}{acender; inflamar; acender uma fogueira, para iluminar}
  \end{Phonetics}
\end{Entry}

%%%%%%%%%% 炼 %%%%%%%%%%
\subsection*{炼}\addcontentsline{loh}{figure}{炼}

\begin{Entry}{炼}{9}{⽕}
  \begin{Phonetics}{炼}{lian4}[][HSK 7-9]
    \definition*{s.}{Sobrenome: Lian}
    \definition{v.}{fundir; refinar; purificar ou endurecer uma substância por meio de aquecimento ou outros métodos | temperar (um metal) com fogo | ponderar as próprias palavras; buscar a expressão adequada; aprimorar; analisar e refinar cuidadosamente as palavras e frases para torná"-las concisas e belas}
  \end{Phonetics}
\end{Entry}

%%%%%%%%%% 烂 %%%%%%%%%%
\subsection*{烂}\addcontentsline{loh}{figure}{烂}

\begin{Entry}{烂}{9}{⽕}
  \begin{Phonetics}{烂}{lan4}[][HSK 5]
    \definition{adj.}{macio; pastoso; amassado | podre; deteriorado | quebrado; esfarrapado; gasto | desorganizado; indigno}
    \definition{adv.}{totalmente; extremamente; completamente; expressa um grau muito profundo}
    \definition{v.}{apodrecer; infeccionar; decompor"-se}
  \end{Phonetics}
\end{Entry}

%%%%%%%%%% 烈 %%%%%%%%%%
\subsection*{烈}\addcontentsline{loh}{figure}{烈}

\begin{Entry}{烈}{10}{⽕}
  \begin{Phonetics}{烈}{lie4}
    \definition*{s.}{Sobrenome: Lie}
    \definition{adj.}{forte; violento; intenso; feroz | justo; severo | firme; convicto; rigoroso}
    \definition{s.}{pessoa que morreu por uma causa justa | conquistas; façanhas | mártir sacrificando"-se por uma causa justa}
  \end{Phonetics}
\end{Entry}

\begin{Entry}{烈士}{10,3}{⽕,⼠}
  \begin{Phonetics}{烈士}{lie4shi4}[][HSK 7-9]
    \definition[位,名]{s.}{mártir; pessoas que se sacrificaram pela causa da justiça | uma pessoa de grande empenho; na antiguidade, referia"-se a uma pessoa que aspirava a alcançar grandes feitos}
  \end{Phonetics}
\end{Entry}

%%%%%%%%%% 烘 %%%%%%%%%%
\subsection*{烘}\addcontentsline{loh}{figure}{烘}

\begin{Entry}{烘}{10}{⽕}
  \begin{Phonetics}{烘}{hong1}
    \definition{v.}{secar; assar; aquecer; usar fogo ou vapor para aquecer o corpo ou para cozinhar, aquecer ou secar algo | destacar}
  \end{Phonetics}
\end{Entry}

\begin{Entry}{烘干}{10,3}{⽕,⼲}
  \begin{Phonetics}{烘干}{hong1gan1}[][HSK 7-9]
    \definition{v.}{secar em fogo alto | secar ao lado ou sobre o fogo | assar; secar no forno}
  \end{Phonetics}
\end{Entry}

\begin{Entry}{烘托}{10,6}{⽕,⼿}
  \begin{Phonetics}{烘托}{hong1tuo1}[][HSK 7-9]
    \definition{v.}{adicionar sombreamento ao redor de um objeto para destacá"-lo; um dos métodos de pintura chinesa, que utiliza tinta ou cores claras para pontilhar o contorno do objeto e torná"-lo mais claro | destacar por contraste; colocar em nítido relevo; fazer com que se destaque}
  \end{Phonetics}
\end{Entry}

%%%%%%%%%% 烟 %%%%%%%%%%
\subsection*{烟}\addcontentsline{loh}{figure}{烟}

\begin{Entry}{烟}{10}{⽕}
  \begin{Phonetics}{烟}{yan1}[][HSK 3]
    \definition[股,支,根,盒,包]{s.}{fumaça; gás produzido pela combustão de materiais, misturado com pequenas partículas não completamente queimadas | névoa; neblina | tabaco; planta de tabaco | fumo; cigarro; termo geral para cigarros, charutos, etc. | ópio | fuligem; fumaça de carvão}
    \definition{v.}{ficar irritado com a fumaça (os olhos lacrimejam ou não conseguem abrir)}
  \end{Phonetics}
\end{Entry}

\begin{Entry}{烟火}{10,4}{⽕,⽕}
  \begin{Phonetics}{烟火}{yan1huo3}[][HSK 7-9]
    \definition{s.}{fumaça e fogo | comida cozida | descendentes; incenso oferecido aos ancestrais; metaforicamente referindo"-se aos descendentes | guerra; fumaça e chamas da guerra; uma metáfora para a guerra}
  \end{Phonetics}
  \begin{Phonetics}{烟火}{yan1huo5}
    \definition{s.}{fogos de artifício}
  \synonymref{烟花}{yan1hua1}
  \synonymref{焰火}{yan4huo3}
  \end{Phonetics}
\end{Entry}

\begin{Entry}{烟叶}{10,5}{⽕,⼝}
  \begin{Phonetics}{烟叶}{yan1ye4}
    \definition{s.}{folha de tabaco | planta de tabaco}
  \end{Phonetics}
\end{Entry}

\begin{Entry}{烟头}{10,5}{⽕,⼤}
  \begin{Phonetics}{烟头}{yan1tou2}
    \definition[根]{s.}{bituca de cigarro}
  \seealsoref{烟屁股}{yan1pi4gu5}
  \seealsoref{烟头儿}{yan1tou2r5}
  \end{Phonetics}
\end{Entry}

\begin{Entry}{烟头儿}{10,5,2}{⽕,⼤,⼉}
  \begin{Phonetics}{烟头儿}{yan1tou2r5}
    \definition{s.}{bituca de cigarro}
  \end{Phonetics}
\end{Entry}

\begin{Entry}{烟囱}{10,7}{⽕,⼞}
  \begin{Phonetics}{烟囱}{yan1cong1}[][HSK 7-9]
    \definition[根]{s.}{chaminé; uma estrutura que proporciona ventilação para gases de combustão quentes ou fumaça provenientes de uma caldeira, fogão, forno ou lareira}
  \synonymref{突}{tu1}
  \end{Phonetics}
\end{Entry}

\begin{Entry}{烟屁股}{10,7,8}{⽕,⼫,⾁}
  \begin{Phonetics}{烟屁股}{yan1pi4gu5}
    \definition{s.}{bituca de cigarro}
  \end{Phonetics}
\end{Entry}

\begin{Entry}{烟花}{10,7}{⽕,⾋}
  \begin{Phonetics}{烟花}{yan1hua1}[][HSK 6]
    \definition[场,朵]{s.}{fogos de artifício; uma coisa que emite faíscas de várias cores quando exposta à observação | prostituta; antigamente, referia"-se a algo relacionado à prostituição}
  \end{Phonetics}
\end{Entry}

\begin{Entry}{烟雨}{10,8}{⽕,⾬}
  \begin{Phonetics}{烟雨}{yan1yu3}
    \definition{s.}{chuva fina; chuvisco; garoa; uma garoa fina como névoa}
  \end{Phonetics}
\end{Entry}

\begin{Entry}{烟草}{10,9}{⽕,⾋}
  \begin{Phonetics}{烟草}{yan1cao3}
    \definition{s.}{tabaco; planta cujas folhas são a principal matéria"-prima para a produção de tabaco}
  \synonymref{香烟}{xiang1yan1}
  \end{Phonetics}
\end{Entry}

%%%%%%%%%% 烤 %%%%%%%%%%
\subsection*{烤}\addcontentsline{loh}{figure}{烤}

\begin{Entry}{烤}{10}{⽕}
  \begin{Phonetics}{烤}{kao3}
    \definition{v.}{assar; tostar; grelhar; cozinhar ou secar um objeto aproximando"-o do fogo | mater"-se quente junto a uma fogueira; aquecer"-se colocando o corpo perto de uma fogueira ou de uma área de alta temperatura}
  \synonymref{烘}{hong1}
  \end{Phonetics}
\end{Entry}

\begin{Entry}{烤肉}{10,6}{⽕,⾁}
  \begin{Phonetics}{烤肉}{kao3rou4}[][HSK 5]
    \definition[块,串,片,盘]{s.}{churrasco (literalmente carne assada)}
  \end{Phonetics}
\end{Entry}

\begin{Entry}{烤鸭}{10,10}{⽕,⿃}
  \begin{Phonetics}{烤鸭}{kao3ya1}[][HSK 5]
    \definition[只,盘]{s.}{pato assado; pato recheado e assado em um forno especial após ser abatido}
  \end{Phonetics}
\end{Entry}

%%%%%%%%%% 烦 %%%%%%%%%%
\subsection*{烦}\addcontentsline{loh}{figure}{烦}

\begin{Entry}{烦}{10}{⽕}
  \begin{Phonetics}{烦}{fan2}[][HSK 4]
    \definition{adj.}{redundante e confuso | supérfluo e confuso; muito bagunçado}
    \definition{v.}{aborrecer | irritar; incomodar; estar cansado de; ficar irritado | incomodar; solicitar}
  \end{Phonetics}
\end{Entry}

\begin{Entry}{烦闷}{10,7}{⽕,⾨}
  \begin{Phonetics}{烦闷}{fan2men4}[][HSK 7-9]
    \definition{adj.}{infeliz; deprimido; mal"-humorado | desconfortável}
  \end{Phonetics}
\end{Entry}

\begin{Entry}{烦恼}{10,9}{⽕,⼼}
  \begin{Phonetics}{烦恼}{fan2nao3}[][HSK 7-9]
    \definition{adj.}{irritado; preocupado; incomodado}
    \definition[个,种,些]{s.}{aborrecimento; coisas que te incomodam}
  \end{Phonetics}
\end{Entry}

\begin{Entry}{烦躁}{10,20}{⽕,⾜}
  \begin{Phonetics}{烦躁}{fan2zao4}[][HSK 7-9]
    \definition{adj.}{inquieto; agitado; irritável}
  \end{Phonetics}
\end{Entry}

%%%%%%%%%% 烧 %%%%%%%%%%
\subsection*{烧}\addcontentsline{loh}{figure}{烧}

\begin{Entry}{烧}{10}{⽕}
  \begin{Phonetics}{烧}{shao1}[][HSK 4]
    \definition[次]{s.}{febre; temperatura corporal mais alta do que o normal}
    \definition{v.}{queimar; pegar fogo | cozinhar; aquecer; assar | guisar depois de fritar ou fritar depois de guisar | assar; grelhar os ingredientes dos alimentos diretamente sobre o fogo | ter febre; estar com febre | danificar (matar ou murchar) as plantas pelo uso excessivo (ou inadequado) de fertilizantes | tornar"-se arrogante ou presunçoso; metáfora de estar em uma boa posição e se deixar levar}
  \end{Phonetics}
\end{Entry}

\begin{Entry}{烧烤}{10,10}{⽕,⽕}
  \begin{Phonetics}{烧烤}{shao1kao3}[][HSK 7-9]
    \definition[顿,次,场]{s.}{churrasco; carne assada ou grelhada}
    \definition{v.}{assar; grelhar; fazer churrasco; grelhar ou assar (carnes) sobre carvão}
  \end{Phonetics}
\end{Entry}

\begin{Entry}{烧毁}{10,13}{⽕,⽎}
  \begin{Phonetics}{烧毁}{shao1hui3}[][HSK 7-9]
    \definition{v.}{queimar; destruir pelo fogo}
  \end{Phonetics}
\end{Entry}

%%%%%%%%%% 烫 %%%%%%%%%%
\subsection*{烫}\addcontentsline{loh}{figure}{烫}

\begin{Entry}{烫}{10}{⽕}
  \begin{Phonetics}{烫}{tang4}[][HSK 7-9]
    \definition{adj.}{escaldante; muito quente; alta temperatura do objeto}
    \definition{v.}{queimar; escaldar; objetos em alta temperatura causam dor ao entrarem em contato com a pele | passar a ferro; prensar; aquecer; aquecer em água quente; utilizar um objeto mais quente para aumentar a temperatura de outro objeto ou provocar outras alterações | fazer permanente (cabelo)}
  \end{Phonetics}
\end{Entry}

%%%%%%%%%% 热 %%%%%%%%%%
\subsection*{热}\addcontentsline{loh}{figure}{热}

\begin{Entry}{热}{10}{⽕}
  \begin{Phonetics}{热}{re4}[][HSK 1]
    \definition{adj.}{quente; temperatura elevada | ardente; caloroso; profundamente afetuoso | ansioso; invejoso; descreve inveja e desejo de possuir algo | térmico; altamente radioativo | popular; muito procurado; muito apreciado por muitas pessoas}
    \definition{s.}{calor; energia liberada pelo movimento irregular das moléculas dentro de um objeto | febre; febre alta causada por doença | moda passageira; mania; febre}
    \definition{v.}{aquecer (geralmente se refere a alimentos)}
  \synonymref{繁}{fan2}
  \synonymref{暖}{nuan3}
  \synonymref{炎}{yan2}
  \antonymref{寒}{han2}
  \antonymref{冷}{leng3}
  \antonymref{凉}{liang2}
  \end{Phonetics}
\end{Entry}

\begin{Entry}{热门}{10,3}{⽕,⾨}
  \begin{Phonetics}{热门}{re4men2}[][HSK 5]
    \definition{adj.}{popular; durante um período de tempo, foi algo que interessava a todos}
    \definition{s.}{algo que desperta o interesse popular; metáfora para algo que está na moda e recebe a atenção de todos (em contraste com 冷门)}
  \seealsoref{冷门}{leng3men2}
  \end{Phonetics}
\end{Entry}

\begin{Entry}{热心}{10,4}{⽕,⼼}
  \begin{Phonetics}{热心}{re4xin1}[][HSK 4]
    \definition{adj.}{ardente; sincero; entusiasmado; afetuoso; apaixonado; interessado}
    \definition{v.}{ser entusiasmado com alguma coisa}
  \end{Phonetics}
\end{Entry}

\begin{Entry}{热气}{10,4}{⽕,⽓}
  \begin{Phonetics}{热气}{re4qi4}[][HSK 7-9]
    \definition[股]{s.}{vapor; calor | Coloquial: entusiasmo | gás quente}
  \end{Phonetics}
\end{Entry}

\begin{Entry}{热气球}{10,4,11}{⽕,⽓,⽟}
  \begin{Phonetics}{热气球}{re4qi4qiu2}[][HSK 7-9]
    \definition{s.}{balão de ar quente}[我们坐了热气球升空。===Fizemos um passeio de balão de ar quente.]
  \end{Phonetics}
\end{Entry}

\begin{Entry}{热水}{10,4}{⽕,⽔}
  \begin{Phonetics}{热水}{re4shui3}[][HSK 6]
    \definition{s.}{água quente; água em temperatura mais alta}
  \end{Phonetics}
\end{Entry}

\begin{Entry}{热水器}{10,4,16}{⽕,⽔,⼝}
  \begin{Phonetics}{热水器}{re4shui3qi4}[][HSK 6]
    \definition[台]{s.}{aquecedor de água; aparelhos que aquecem água usando eletricidade, gás natural, gás liquefeito de petróleo ou energia solar}
  \end{Phonetics}
\end{Entry}

\begin{Entry}{热血沸腾}{10,6,8,13}{⽕,⾎,⽔,⾁}
  \begin{Phonetics}{热血沸腾}{re4xue4-fei4teng2}
    \definition{expr.}{estar animado; ter o sangue correndo}
  \end{Phonetics}
\end{Entry}

\begin{Entry}{热泪盈眶}{10,8,9,11}{⽕,⽔,⽫,⽬}
  \begin{Phonetics}{热泪盈眶}{re4lei4ying2kuang4}
    \definition{expr.}{com os olhos cheios de lágrimas; extremamente comovido; olhos cheios de lágrimas, expressando extrema gratidão ou alegria}
  \antonymref{眉开眼笑}{mei2kai1-yan3xiao4}
  \end{Phonetics}
\end{Entry}

\begin{Entry}{热线}{10,8}{⽕,⽷}
  \begin{Phonetics}{热线}{re4xian4}[][HSK 6]
    \definition[条]{s.}{raio infravermelho | linha direta; \emph{hot line}; uma linha telefônica ou telegráfica direta; uma linha para um ponto de acesso | rota quente (ou movimentada, popular) | raio de calor}
  \end{Phonetics}
\end{Entry}

\begin{Entry}{热闹}{10,8}{⽕,⾾}
  \begin{Phonetics}{热闹}{re4nao5}[][HSK 4]
    \definition{adj.}{animado; agitado; movimentado com barulho e excitação; descreve uma cena animada com uma atmosfera calorosa}
    \definition{s.}{uma vista emocionante; uma cena de agitação e excitação; atmosfera acolhedora}
    \definition{v.}{animar; divertir"-se}
  \end{Phonetics}
\end{Entry}

\begin{Entry}{热带}{10,9}{⽕,⼱}
  \begin{Phonetics}{热带}{re4dai4}[][HSK 7-9]
    \definition{s.}{zona tropical; os trópicos; localizada em ambos os lados da linha do Equador, entre o Trópico de Câncer e o Trópico de Capricórnio, a região apresenta condições climáticas onde a duração do dia e da noite varia pouco com as estações do ano, o clima é quente durante todo o ano e as chuvas são abundantes}
  \end{Phonetics}
\end{Entry}

\begin{Entry}{热点}{10,9}{⽕,⽕}
  \begin{Phonetics}{热点}{re4dian3}[][HSK 6]
    \definition{s.}{ponto de acesso; \emph{hotspot}}
  \end{Phonetics}
\end{Entry}

\begin{Entry}{热烈}{10,10}{⽕,⽕}
  \begin{Phonetics}{热烈}{re4lie4}[][HSK 3]
    \definition{adj.}{caloroso; fervoroso; ardente; entusiasmado; excitado}
  \end{Phonetics}
\end{Entry}

\begin{Entry}{热爱}{10,10}{⽕,⽖}
  \begin{Phonetics}{热爱}{re4'ai4}[][HSK 3]
    \definition{v.}{amar ardentemente; amar de coração; ter amor profundo por; amar apaixonadamente}
  \end{Phonetics}
\end{Entry}

\begin{Entry}{热衷}{10,10}{⽕,⾐}
  \begin{Phonetics}{热衷}{re4zhong1}[][HSK 7-9]
    \definition{v.}{gostar de; ter grande apreço por; gostar muito de (uma determinada atividade) | ansiar; desejar ardentemente; buscar avidamente por (fama, fortuna e poder)}
  \end{Phonetics}
\end{Entry}

\begin{Entry}{热情}{10,11}{⽕,⼼}
  \begin{Phonetics}{热情}{re4qing2}[][HSK 2]
    \definition{adj.}{caloroso; fervoroso; entusiasmado; cordial; descreve sentimentos calorosos por alguém}
    \definition{s.}{entusiasmo; ardor; devoção; calor humano; zelo; sentimentos calorosos}
  \seealsoref{热烈}{re4lie4}
  \seealsoref{热心}{re4xin1}
  \synonymref{激情}{ji1qing2}
  \synonymref{亲切}{qin1qie4}
  \synonymref{亲热}{qin1re4}
  \synonymref{热烈}{re4lie4}
  \synonymref{热心}{re4xin1}
  \synonymref{殷勤}{yin1qin2}
  \antonymref{冷淡}{leng3dan4}
  \antonymref{冷酷}{leng3ku4}
  \antonymref{冷落}{leng3luo4}
  \antonymref{冷漠}{leng3mo4}
  \end{Phonetics}
\end{Entry}

\begin{Entry}{热量}{10,12}{⽕,⾥}
  \begin{Phonetics}{热量}{re4liang4}[][HSK 5]
    \definition{s.}{calor; quantidade de calor; calorias; em física, refere"-se à energia transferida entre objetos com temperaturas diferentes, do objeto com temperatura mais alta para o objeto com temperatura mais baixa}
  \end{Phonetics}
\end{Entry}

\begin{Entry}{热腾腾}{10,13,13}{⽕,⾁,⾁}
  \begin{Phonetics}{热腾腾}{re4teng2teng2}[][HSK 7-9]
    \definition{adj.}{bem quente; fumegando}[桌上放着一碗热腾腾的汤。===Uma tigela de sopa fumegante está sobre a mesa.]
  \end{Phonetics}
\end{Entry}

\begin{Entry}{热潮}{10,15}{⽕,⽔}
  \begin{Phonetics}{热潮}{re4chao2}[][HSK 7-9]
    \definition{s.}{surto; mania; onda de entusiasmo; descreve uma situação próspera e agitada}
  \end{Phonetics}
\end{Entry}

%%%%%%%%%% 惨 %%%%%%%%%%
\subsection*{惨}\addcontentsline{loh}{figure}{惨}

\begin{Entry}{惨}{11}{⽕}
  \begin{Phonetics}{惨}{can3}[][HSK 6]
    \definition{adj.}{miserável; trágico | cruel; brutal; implacável | desastroso; terrível; esmagador | lamentável; desaventurado | em um grau sério; grau grave; dano grave | selvagem; desumano; vicioso; cruel}
  \end{Phonetics}
\end{Entry}

\begin{Entry}{惨白}{11,5}{⽕,⽩}
  \begin{Phonetics}{惨白}{can3bai2}[][HSK 7-9]
    \definition{adj.}{(cenário) escuro; fraco; sombrio; sombrio | (rosto) mortalmente pálido; medonho | terrivelmente (fantasmagórico) pálido; pálido; escuro}
  \end{Phonetics}
\end{Entry}

\begin{Entry}{惨重}{11,9}{⽕,⾥}
  \begin{Phonetics}{惨重}{can3zhong4}[][HSK 7-9]
    \definition{adj.}{pesado; doloroso; desastroso | calamitoso (perdas extremamente graves)}
  \end{Phonetics}
\end{Entry}

\begin{Entry}{惨痛}{11,12}{⽕,⽧}
  \begin{Phonetics}{惨痛}{can3tong4}[][HSK 7-9]
    \definition{adj.}{grave; terrivelmente doloroso; amargo; agonizante | profundamente triste; doloroso}
  \end{Phonetics}
\end{Entry}

%%%%%%%%%% 烹 %%%%%%%%%%
\subsection*{烹}\addcontentsline{loh}{figure}{烹}

\begin{Entry}{烹}{11}{⽕}
  \begin{Phonetics}{烹}{peng1}
    \definition{v.}{ferver; cozinhar | saltear; refogar; fritar}
  \end{Phonetics}
\end{Entry}

\begin{Entry}{烹调}{11,10}{⽕,⾔}
  \begin{Phonetics}{烹调}{peng1tiao2}[][HSK 7-9]
    \definition{v.}{cozinhar; cozinhar e temperar}
  \end{Phonetics}
\end{Entry}

%%%%%%%%%% 焊 %%%%%%%%%%
\subsection*{焊}\addcontentsline{loh}{figure}{焊}

\begin{Entry}{焊}{11}{⽕}
  \begin{Phonetics}{焊}{han4}[][HSK 7-9]
    \definition{v.}{soldar; usar metal fundido para reparar objetos de metal ou conectar peças de metal}
  \end{Phonetics}
\end{Entry}

%%%%%%%%%% 焕 %%%%%%%%%%
\subsection*{焕}\addcontentsline{loh}{figure}{焕}

\begin{Entry}{焕}{11}{⽕}
  \begin{Phonetics}{焕}{huan4}
    \definition{adj.}{Literário: brilhante; reluzente; radiante}
  \end{Phonetics}
\end{Entry}

\begin{Entry}{焕发}{11,5}{⽕,⼜}
  \begin{Phonetics}{焕发}{huan4fa1}[][HSK 7-9]
    \definition{v.}{brilhar; revigorar; irradiar}
  \end{Phonetics}
\end{Entry}

%%%%%%%%%% 悲 %%%%%%%%%%
\subsection*{悲}\addcontentsline{loh}{figure}{悲}

\begin{Entry}{悲}{12}{⽕}
  \begin{Phonetics}{悲}{bei1}
    \definition{adj.}{triste; pesaroso; melancólico | compassivo; misericordioso}
  \synonymref{哀}{ai1}
  \antonymref{欢}{huan1}
  \antonymref{乐}{le4}
  \antonymref{喜}{xi3}
  \end{Phonetics}
\end{Entry}

\begin{Entry}{悲伤}{12,6}{⽕,⼈}
  \begin{Phonetics}{悲伤}{bei1shang1}[][HSK 5]
    \definition{adj.}{triste; pesaroso}
  \seealsoref{悲哀}{bei1'ai1}
  \seealsoref{悲痛}{bei1tong4}
  \end{Phonetics}
\end{Entry}

\begin{Entry}{悲欢离合}{12,6,10,6}{⽕,⽋,⼇,⼝}
  \begin{Phonetics}{悲欢离合}{bei1huan1-li2he2}[][HSK 7-9]
    \definition{expr.}{alegrias e tristezas | separações e reencontros | as vicissitudes da vida}
  \end{Phonetics}
\end{Entry}

\begin{Entry}{悲观}{12,6}{⽕,⾒}
  \begin{Phonetics}{悲观}{bei1guan1}[][HSK 7-9]
    \definition{adj.}{pessimista; negativismo, falta de confiança no futuro}
  \synonymref{悲哀}{bei1'ai1}
  \synonymref{灰心}{hui1/xin1}
  \synonymref{绝望}{jue2/wang4}
  \synonymref{扫兴}{sao3/xing4}
  \synonymref{失望}{shi1wang4}
  \synonymref{颓废}{tui2fei4}
  \synonymref{消沉}{xiao1chen2}
  \synonymref{消极}{xiao1ji2}
  \antonymref{乐观}{le4guan1}
  \end{Phonetics}
\end{Entry}

\begin{Entry}{悲哀}{12,9}{⽕,⼝}
  \begin{Phonetics}{悲哀}{bei1'ai1}[][HSK 7-9]
    \definition{adj.}{triste}
    \definition{s.}{tristeza; refere"-se a coisas tristes e dolorosas}
  \seealsoref{悲伤}{bei1shang1}
  \synonymref{悲痛}{bei1tong4}
  \antonymref{欢乐}{huan1le4}
  \end{Phonetics}
\end{Entry}

\begin{Entry}{悲剧}{12,10}{⽕,⼑}
  \begin{Phonetics}{悲剧}{bei1ju4}[][HSK 5]
    \definition[出,部]{s.}{tragédia; drama trágico; uma das principais categorias de teatro, caracterizada basicamente pela representação do conflito irreconciliável entre o protagonista e a realidade e seu final trágico | tragédia; evento triste; metáfora para encontro infeliz}
  \antonymref{喜剧}{xi3ju4}
  \end{Phonetics}
\end{Entry}

\begin{Entry}{悲惨}{12,11}{⽕,⽕}
  \begin{Phonetics}{悲惨}{bei1can3}[][HSK 6]
    \definition{adj.}{trágico; miserável; extremamente doloroso e triste}
  \seealsoref{凄惨}{qi1can3}
  \synonymref{悲哀}{bei1'ai1}
  \synonymref{凄惨}{qi1can3}
  \synonymref{凄凉}{qi1liang2}
  \synonymref{痛苦}{tong4ku3}
  \antonymref{幸福}{xing4fu2}
  \end{Phonetics}
\end{Entry}

\begin{Entry}{悲痛}{12,12}{⽕,⽧}
  \begin{Phonetics}{悲痛}{bei1tong4}[][HSK 7-9]
    \definition{adj.}{triste; aflito}
    \definition{s.}{tristeza; sofrimento}
  \seealsoref{悲伤}{bei1shang1}
  \synonymref{悲哀}{bei1'ai1}
  \synonymref{悲伤}{bei1shang1}
  \synonymref{沮丧}{ju3sang4}
  \synonymref{伤心}{shang1/xin1}
  \antonymref{高兴}{gao1xing4}
  \antonymref{开心}{kai1/xin1}
  \antonymref{快乐}{kuai4le4}
  \antonymref{喜悦}{xi3yue4}
  \antonymref{愉快}{yu2kuai4}
  \end{Phonetics}
\end{Entry}

%%%%%%%%%% 焚 %%%%%%%%%%
\subsection*{焚}\addcontentsline{loh}{figure}{焚}

\begin{Entry}{焚}{12}{⽕}
  \begin{Phonetics}{焚}{fen2}
    \definition{v.}{queimar}
  \end{Phonetics}
\end{Entry}

\begin{Entry}{焚香}{12,9}{⽕,⾹}
  \begin{Phonetics}{焚香}{fen2xiang1}
    \definition{v.}{queimar incenso}
  \end{Phonetics}
\end{Entry}

\begin{Entry}{焚烧}{12,10}{⽕,⽕}
  \begin{Phonetics}{焚烧}{fen2shao1}[][HSK 7-9]
    \definition{v.}{queimar; incendiar; colocar fogo}
  \end{Phonetics}
\end{Entry}

%%%%%%%%%% 焦 %%%%%%%%%%
\subsection*{焦}\addcontentsline{loh}{figure}{焦}

\begin{Entry}{焦}{12}{⽕}
  \begin{Phonetics}{焦}{jiao1}[][HSK 7-9]
    \definition*{s.}{Sobrenome: Jiao}
    \definition{adj.}{queimado; chamuscado; carbonizado | preocupado; ansioso}
    \definition{clas.}{J; Joule, abreviação}
    \definition{pref.}{(química) piro-}
    \definition{s.}{Metalurgia: coque}
  \end{Phonetics}
\end{Entry}

\begin{Entry}{焦急}{12,9}{⽕,⼼}
  \begin{Phonetics}{焦急}{jiao1ji2}[][HSK 7-9]
    \definition{adj.}{ansioso; preocupado; com pressa}
  \end{Phonetics}
\end{Entry}

\begin{Entry}{焦点}{12,9}{⽕,⽕}
  \begin{Phonetics}{焦点}{jiao1dian3}[][HSK 6]
    \definition{s.}{foco; ponto focal; Matemática: refere"-se a um ponto que tem uma relação especial com uma elipse, hipérbole, parábola, etc. | foco; ponto focal; Óptica: refere"-se à intersecção de feixes de luz paralelos após serem refratados por uma lente ou refletidos por um espelho curvo | foco; questão central; metaforicamente, uma coisa ou princípio que chama a atenção para o foco}
  \end{Phonetics}
\end{Entry}

\begin{Entry}{焦虑}{12,10}{⽕,⾌}
  \begin{Phonetics}{焦虑}{jiao1lv4}[][HSK 7-9]
    \definition{adj.}{ansioso; preocupado; apreensivo}
  \end{Phonetics}
\end{Entry}

\begin{Entry}{焦距}{12,11}{⽕,⾜}
  \begin{Phonetics}{焦距}{jiao1ju4}[][HSK 7-9]
    \definition{s.}{Ótica: distância focal; comprimento focal}
  \end{Phonetics}
\end{Entry}

\begin{Entry}{焦躁}{12,20}{⽕,⾜}
  \begin{Phonetics}{焦躁}{jiao1zao4}[][HSK 7-9]
    \definition{adj.}{impaciente; inquieto de ansiedade; ansioso e irritadiço}
  \end{Phonetics}
\end{Entry}

%%%%%%%%%% 焰 %%%%%%%%%%
\subsection*{焰}\addcontentsline{loh}{figure}{焰}

\begin{Entry}{焰}{12}{⽕}
  \begin{Phonetics}{焰}{yan4}
    \definition{adj.}{arrogante; metáfora para poder e prestígio}
    \definition{s.}{chama; labareda}
  \end{Phonetics}
\end{Entry}

\begin{Entry}{焰火}{12,4}{⽕,⽕}
  \begin{Phonetics}{焰火}{yan4huo3}[][HSK 7-9]
    \definition[种]{s.}{fogos de artifício}
  \synonymref{火花}{huo3hua1}
  \synonymref{烟火}{yan1huo5}
  \end{Phonetics}
\end{Entry}

%%%%%%%%%% 然 %%%%%%%%%%
\subsection*{然}\addcontentsline{loh}{figure}{然}

\begin{Entry}{然}{12}{⽕}
  \begin{Phonetics}{然}{ran2}
    \definition*{s.}{Sobrenome: Ran}
    \definition{adj.}{certo; correto}
    \definition{adv.}{então; assim; desta forma; assim mesmo}
    \definition{conj.}{mas; contudo; porém; no entanto}
    \definition{suf.}{sufixos adverbiais ou adjetivais}
  \end{Phonetics}
\end{Entry}

\begin{Entry}{然后}{12,6}{⽕,⼝}
  \begin{Phonetics}{然后}{ran2hou4}[][HSK 2]
    \definition{conj.}{então; depois disso; posteriormente; indica que algo segue após uma ação ou situação}
  \seealsoref{后来}{hou4lai2}
  \seealsoref{以后}{yi3hou4}
  \synonymref{后来}{hou4lai2}
  \synonymref{继而}{ji4'er2}
  \synonymref{接着}{jie1zhe5}
  \synonymref{其次}{qi2ci4}
  \synonymref{以后}{yi3hou4}
  \synonymref{之后}{zhi1hou4}
  \end{Phonetics}
\end{Entry}

\begin{Entry}{然而}{12,6}{⽕,⽽}
  \begin{Phonetics}{然而}{ran2'er2}[][HSK 4]
    \definition{conj.}{ainda; mas; contudo; todavia; usado no início de uma frase para indicar uma transição; para indicar uma transição, geralmente é precedido por uma conjunção como 虽然 para indicar concessão}
  \seealsoref{虽然}{sui1ran2}
  \end{Phonetics}
\end{Entry}

%%%%%%%%%% 煮 %%%%%%%%%%
\subsection*{煮}\addcontentsline{loh}{figure}{煮}

\begin{Entry}{煮}{12}{⽕}
  \begin{Phonetics}{煮}{zhu3}[][HSK 6]
    \definition*{s.}{Sobrenome: Zhu}
    \definition{v.}{ferver; cozinhar; aquecer alimentos ou outros itens em água}
  \end{Phonetics}
\end{Entry}

%%%%%%%%%% 煎 %%%%%%%%%%
\subsection*{煎}\addcontentsline{loh}{figure}{煎}

\begin{Entry}{煎}{13}{⽕}
  \begin{Phonetics}{煎}{jian1}[][HSK 7-9]
    \definition{s.}{Fitoterapia: decocção}
    \definition{v.}{fritar (em pouco óleo, sem mexer) | ferver (em água)}
  \end{Phonetics}
\end{Entry}

\begin{Entry}{煎饼}{13,9}{⽕,⾷}
  \begin{Phonetics}{煎饼}{jian1bing3}
    \definition[张,块,摞]{s.}{uma panqueca fina feita de farinha de milho"-miúdo}
  \end{Phonetics}
\end{Entry}

\begin{Entry}{煎蛋}{13,11}{⽕,⾍}
  \begin{Phonetics}{煎蛋}{jian1dan4}
    \definition{s.}{ovo frito}
  \end{Phonetics}
\end{Entry}

%%%%%%%%%% 煞 %%%%%%%%%%
\subsection*{煞}\addcontentsline{loh}{figure}{煞}

\begin{Entry}{煞}{13}{⽕}
  \begin{Phonetics}{煞}{sha1}
    \definition{v.}{parar; terminar; cessar; trancar; concluir; finalizar | frear; conter; apertar; prender | diminuir gradualmente; enfraquecer; prejudicar; danificar}
  \end{Phonetics}
  \begin{Phonetics}{煞}{sha4}
    \definition{adv.}{muito; extremamente}
    \definition{s.}{\emph{goblin}; duende; espírito maligno; deus maligno}
    \definition{v.}{parar; interromper; verificar; encerrar | apertar}
  \end{Phonetics}
\end{Entry}

\begin{Entry}{煞气}{13,4}{⽕,⽓}
  \begin{Phonetics}{煞气}{sha1qi4}
    \definition{v.}{desabafar | descontar em alguém | descarregar a raiva em (uma pessoa inocente)}
  \end{Phonetics}
  \begin{Phonetics}{煞气}{sha4qi4}
    \definition{s.}{olhar sinistro | influência desfavorável}
    \definition{v.}{(pneu, balão, etc.) vazar ar}
  \end{Phonetics}
\end{Entry}

%%%%%%%%%% 煤 %%%%%%%%%%
\subsection*{煤}\addcontentsline{loh}{figure}{煤}

\begin{Entry}{煤}{13}{⽕}
  \begin{Phonetics}{煤}{mei2}[][HSK 5]
    \definition[块,吨,斤,堆]{s.}{carvão; carvão vegetal; minério sólido preto}
  \end{Phonetics}
\end{Entry}

\begin{Entry}{煤气}{13,4}{⽕,⽓}
  \begin{Phonetics}{煤气}{mei2qi4}[][HSK 5]
    \definition[罐,瓶]{s.}{gás; gás de carvão; gás obtido a partir do processamento do carvão não tem cor nem odor, é tóxico e pode ser queimado ou utilizado como matéria"-prima na indústria química | envenenamento por monóxido de carbono}
  \end{Phonetics}
\end{Entry}

\begin{Entry}{煤矿}{13,8}{⽕,⽯}
  \begin{Phonetics}{煤矿}{mei2kuang4}[][HSK 7-9]
    \definition{s.}{mina de carvão}
  \end{Phonetics}
\end{Entry}

\begin{Entry}{煤炭}{13,9}{⽕,⽕}
  \begin{Phonetics}{煤炭}{mei2tan4}[][HSK 7-9]
    \definition{s.}{carvão}
  \end{Phonetics}
\end{Entry}

%%%%%%%%%% 照 %%%%%%%%%%
\subsection*{照}\addcontentsline{loh}{figure}{照}

\begin{Entry}{照}{13}{⽕}
  \begin{Phonetics}{照}{zhao4}[][HSK 3]
    \definition{adv.}{de acordo com; significa agir de acordo com o original ou com determinados padrões}
    \definition{prep.}{em direção a; na direção de | de acordo com; em conformidade com}
    \definition{s.}{imagem; fotografia | permissão; licença; autorização | brilho; iluminação}
    \definition{v.}{brilhar; iluminar | refletir; espelhar; olhar para sua própria imagem em um espelho, etc. | filmar; fotografar; tirar uma foto (fotografia) | cuidar de; tomar conta de; zelar por | notificar; informar | contrastar; comparar; verificar | entender; compreender}
  \end{Phonetics}
\end{Entry}

\begin{Entry}{照片}{13,4}{⽕,⽚}
  \begin{Phonetics}{照片}{zhao4pian4}[][HSK 2]
    \definition[张,套,幅]{s.}{fotografia, foto, imagem}
  \synonymref{图片}{tu2pian4}
  \synonymref{相片}{xiang4pian4}
  \end{Phonetics}
\end{Entry}

\begin{Entry}{照片子}{13,4,3}{⽕,⽚,⼦}
  \begin{Phonetics}{照片子}{zhao4pian4 zi5}
    \definition{v.}{tirar um raio X}
  \end{Phonetics}
\end{Entry}

\begin{Entry}{照片底版}{13,4,8,8}{⽕,⽚,⼴,⽚}
  \begin{Phonetics}{照片底版}{zhao4pian4 di3ban3}
    \definition{s.}{placa fotográfica; negativo fotográfico}
  \end{Phonetics}
\end{Entry}

\begin{Entry}{照应}{13,7}{⽕,⼴}
  \begin{Phonetics}{照应}{zhao4ying4}
    \definition{v.}{coordenar; correlacionar; responder}
  \end{Phonetics}
  \begin{Phonetics}{照应}{zhao4ying5}
    \definition{v.}{cuidar de; zelar por}
  \seealsoref{照料}{zhao4liao4}
  \synonymref{顾问}{gu4wen4}
  \synonymref{关照}{guan1zhao4}
  \synonymref{呼应}{hu1ying4}
  \synonymref{看护}{kan1hu4}
  \synonymref{照料}{zhao4liao4}
  \synonymref{照顾}{zhao4gu5}
  \antonymref{欺负}{qi1fu5}
  \end{Phonetics}
\end{Entry}

\begin{Entry}{照例}{13,8}{⽕,⼈}
  \begin{Phonetics}{照例}{zhao4li4}[][HSK 7-9]
    \definition{adv.}{geralmente; como de costume; em regra; de acordo com o costume e o bom senso}
  \seealsoref{照样}{zhao4yang4}
  \antonymref{反常}{fan3chang2}
  \antonymref{异常}{yi4chang2}
  \end{Phonetics}
\end{Entry}

\begin{Entry}{照明}{13,8}{⽕,⽇}
  \begin{Phonetics}{照明}{zhao4ming2}[][HSK 7-9]
    \definition{v.}{iluminar; iluminar com uma lâmpada ou outra fonte de luz}
  \end{Phonetics}
\end{Entry}

\begin{Entry}{照亮}{13,9}{⽕,⼇}
  \begin{Phonetics}{照亮}{zhao4liang4}
    \definition{s.}{iluminação}
    \definition{v.}{iluminar}
  \end{Phonetics}
\end{Entry}

\begin{Entry}{照相}{13,9}{⽕,⽬}
  \begin{Phonetics}{照相}{zhao4/xiang4}[][HSK 2]
    \definition{v.+compl.}{fotografar; tirar fotos; tirar uma foto; tirar uma fotografia}
  \synonymref{合影}{he2/ying3}
  \synonymref{拍照}{pai1/zhao4}
  \synonymref{拍摄}{pai1she4}
  \synonymref{摄影}{she4ying3}
  \end{Phonetics}
\end{Entry}

\begin{Entry}{照相机}{13,9,6}{⽕,⽬,⽊}
  \begin{Phonetics}{照相机}{zhao4xiang4ji1}
    \definition[部,个,架,台,只]{s.}{câmera; máquina fotográfica}
  \synonymref{摄像机}{she4xiang4ji1}
  \end{Phonetics}
\end{Entry}

\begin{Entry}{照准}{13,10}{⽕,⼎}
  \begin{Phonetics}{照准}{zhao4zhun3}
    \definition{s.}{solicitação concedida (uso formal em documento antigo)}
    \definition{v.}{mirar (arma)}
  \end{Phonetics}
\end{Entry}

\begin{Entry}{照料}{13,10}{⽕,⽃}
  \begin{Phonetics}{照料}{zhao4liao4}[][HSK 7-9]
    \definition{v.}{atender a; cuidar de}
  \seealsoref{照应}{zhao4ying5}
  \synonymref{办理}{ban4li3}
  \synonymref{处理}{chu3li3}
  \synonymref{顾问}{gu4wen4}
  \synonymref{关照}{guan1zhao4}
  \synonymref{管理}{guan3li3}
  \synonymref{料理}{liao4li3}
  \synonymref{收拾}{shou1shi5}
  \synonymref{照顾}{zhao4gu5}
  \synonymref{照应}{zhao4ying5}
  \antonymref{伤害}{shang1hai4}
  \end{Phonetics}
\end{Entry}

\begin{Entry}{照样}{13,10}{⽕,⽊}
  \begin{Phonetics}{照样}{zhao4yang4}[][HSK 6]
    \definition{adv.}{como antes; da mesma maneira; isso significa que, embora as condições externas tenham mudado ou sido afetadas, uma determinada situação ou estado permanece inalterado | após um padrão; de um modelo}
    \definition{v.}{fazer algo de uma certa maneira}
  \end{Phonetics}
\end{Entry}

\begin{Entry}{照顾}{13,10}{⽕,⾴}
  \begin{Phonetics}{照顾}{zhao4gu5}[][HSK 2]
    \definition{v.}{cuidar; cuidar de; atender | oferecer tratamento preferencial; prestar atenção especial e dar tratamento preferencial | (de um cliente) patrocinar; comprar em; cuidar de clientes que vêm comprar coisas ou solicitar serviços em lojas ou indústrias de serviços | dar consideração a; mostrar consideração por; levar em conta; fazer concessões a}
  \seealsoref{关照}{guan1zhao4}
  \synonymref{关心}{guan1xin1}
  \synonymref{关照}{guan1zhao4}
  \synonymref{看护}{kan1hu4}
  \synonymref{体贴}{ti3tie1}
  \synonymref{照料}{zhao4liao4}
  \synonymref{照应}{zhao4ying5}
  \antonymref{忽视}{hu1shi4}
  \antonymref{冷落}{leng3luo4}
  \antonymref{欺负}{qi1fu5}
  \antonymref{无视}{wu2shi4}
  \end{Phonetics}
\end{Entry}

\begin{Entry}{照常}{13,11}{⽕,⼱}
  \begin{Phonetics}{照常}{zhao4chang2}[][HSK 7-9]
    \definition{adv.}{como de costume; isso indica que a situação permanece inalterada}
  \antonymref{偶尔}{ou3'er3}
  \end{Phonetics}
\end{Entry}

\begin{Entry}{照骗}{13,12}{⽕,⾺}
  \begin{Phonetics}{照骗}{zhao4pian4}
    \definition{s.}{Gíria da \emph{Internet}:  foto lisonjeira (trocadilho com 照片) | imagem alterada digitalmente; ``photoshopada''}
  \seealsoref{照片}{zhao4pian4}
  \end{Phonetics}
\end{Entry}

\begin{Entry}{照像}{13,13}{⽕,⼈}
  \begin{Phonetics}{照像}{zhao4/xiang4}
    \variantof{照相}
  \end{Phonetics}
\end{Entry}

\begin{Entry}{照像机}{13,13,6}{⽕,⼈,⽊}
  \begin{Phonetics}{照像机}{zhao4xiang4ji1}
    \variantof{照相机}
  \end{Phonetics}
\end{Entry}

\begin{Entry}{照搬}{13,13}{⽕,⼿}
  \begin{Phonetics}{照搬}{zhao4/ban1}[][HSK 7-9]
    \definition{v.+compl.}{imitar rigidamente; imitar indiscriminadamente; copiar; seguir de perto}
  \end{Phonetics}
\end{Entry}

\begin{Entry}{照耀}{13,20}{⽕,⽻}
  \begin{Phonetics}{照耀}{zhao4yao4}[][HSK 6]
    \definition{v.}{brilhar; iluminar | esclarecer}
  \end{Phonetics}
\end{Entry}

%%%%%%%%%% 煲 %%%%%%%%%%
\subsection*{煲}\addcontentsline{loh}{figure}{煲}

\begin{Entry}{煲}{13}{⽕}
  \begin{Phonetics}{煲}{bao1}[][HSK 7-9]
    \definition{s.}{panela; caldeira}
    \definition{v.}{cozinhar; ensopar; ferver}
  \end{Phonetics}
\end{Entry}

%%%%%%%%%% 蒸 %%%%%%%%%%
\subsection*{蒸}\addcontentsline{loh}{figure}{蒸}

\begin{Entry}{蒸}{13}{⽕}
  \begin{Phonetics}{蒸}{zheng1}[][HSK 7-9]
    \definition{s.}{Arcaico: tocha feita de talos de cânhamo ou bambu | Arcaico: lenha finamente picada}
    \definition{v.}{evaporar; cozinhar no vapor}
  \end{Phonetics}
\end{Entry}

%%%%%%%%%% 愿 %%%%%%%%%%
\subsection*{愿}\addcontentsline{loh}{figure}{愿}

\begin{Entry}{愿}{14}{⽕}
  \begin{Phonetics}{愿}{yuan4}[][HSK 5]
    \definition{adj.}{honesto e prudente}
    \definition{s.}{esperança; desejo; vontade; a ideia de alcançar algum objetivo no futuro | voto (feito perante o Buda ou um deus); o desejo de retribuição feito ao rezar para os deuses e Buda}
    \definition{v.}{estar disposto; estar pronto; de bom grado, concordar porque está de acordo com seus desejos | ter esperança; desejar; qerer alcançar algum desejo}
  \end{Phonetics}
\end{Entry}

\begin{Entry}{愿望}{14,11}{⽕,⽉}
  \begin{Phonetics}{愿望}{yuan4wang4}[][HSK 3]
    \definition[个,种]{s.}{desejo; aspiração; a ideia de alcançar algum objetivo no futuro.}
  \end{Phonetics}
\end{Entry}

\begin{Entry}{愿意}{14,13}{⽕,⼼}
  \begin{Phonetics}{愿意}{yuan4yi5}[][HSK 2]
    \definition{v.}{estar disposto; estar pronto | desejar; ter esperança}
  \seealsoref{乐趣}{le4qu4}
  \synonymref{承诺}{cheng2nuo4}
  \synonymref{渴望}{ke3wang4}
  \synonymref{乐意}{le4yi4}
  \synonymref{盼望}{pan4wang4}
  \synonymref{期望}{qi1wang4}
  \synonymref{情愿}{qing2yuan4}
  \synonymref{同意}{tong2yi4}
  \synonymref{希望}{xi1wang4}
  \synonymref{允许}{yun3xu3}
  \antonymref{不肯}{bu4ken3}
  \antonymref{反对}{fan3dui4}
  \end{Phonetics}
\end{Entry}

%%%%%%%%%% 煽 %%%%%%%%%%
\subsection*{煽}\addcontentsline{loh}{figure}{煽}

\begin{Entry}{煽}{14}{⽕}
  \begin{Phonetics}{煽}{shan1}
    \definition{v.}{abanar (fogo); agitar um leque ou outra folha | incitar; instigar; agitar | vangloriar"-se de; esbanjar prêmios em}
  \end{Phonetics}
\end{Entry}

\begin{Entry}{煽动}{14,6}{⽕,⼒}
  \begin{Phonetics}{煽动}{shan1dong4}[][HSK 7-9]
    \definition{v.}{instigar; incitar (alguém a fazer coisas ruins); agitar; inflamar}
  \end{Phonetics}
\end{Entry}

%%%%%%%%%% 熄 %%%%%%%%%%
\subsection*{熄}\addcontentsline{loh}{figure}{熄}

\begin{Entry}{熄}{14}{⽕}
  \begin{Phonetics}{熄}{xi1}
    \definition{v.}{extinguir; apagar | (fogo, luz, etc.) apagar"-se; extinguir"-se}
  \end{Phonetics}
\end{Entry}

\begin{Entry}{熄火}{14,4}{⽕,⽕}
  \begin{Phonetics}{熄火}{xi1/huo3}[][HSK 7-9]
    \definition{v.+compl.}{(combustível, fogão, etc.) parar de queimar; extinguir"-se | (um motor) parar de funcionar; morrer | interromper a queima (de combustível); desligar; parar (um motor, etc.) | (competições) acabar; finalizar}
  \antonymref{发动}{fa1dong4}
  \end{Phonetics}
\end{Entry}

%%%%%%%%%% 熊 %%%%%%%%%%
\subsection*{熊}\addcontentsline{loh}{figure}{熊}

\begin{Entry}{熊}{14}{⽕}
  \begin{Phonetics}{熊}{xiong2}[][HSK 5]
    \definition*{s.}{Sobrenome: Xiong}
    \definition[头,只]{s.}{urso}
    \definition{v.}{repreender; censurar}
  \end{Phonetics}
\end{Entry}

\begin{Entry}{熊猫}{14,11}{⽕,⽝}
  \begin{Phonetics}{熊猫}{xiong2mao1}
    \definition[只,头]{s.}{panda gigante}
  \seealsoref{猫熊}{mao1xiong2}
  \end{Phonetics}
\end{Entry}

%%%%%%%%%% 熏 %%%%%%%%%%
\subsection*{熏}\addcontentsline{loh}{figure}{熏}

\begin{Entry}{熏}{14}{⽕}
  \begin{Phonetics}{熏}{xun1}[][HSK 7-9]
    \definition{adj.}{quente; suave; ameno}
    \definition{v.}{fumar; fumigar; expor à fumaça ou vapor; (fumaça, vapores, etc.) que entram em contato com um objeto fazem com que ele mude de cor ou adquira um odor | defumar (alimentos); curar por meio da fumaça; métodos de processamento de alimentos que utilizam fumaça e fogo para grelhar os alimentos, conferindo"-lhes um sabor especial | permear; influenciar; estar exposto a; afetado por exposição prolongada}
  \end{Phonetics}
  \begin{Phonetics}{熏}{xun4}
    \definition{adj.}{Literário: quente; ameno}
    \definition{v.}{expor à fumaça ou vapores; fumigar | tratar (carnes, peixes, etc.) com fumaça; defumar | perfumar com incenso, etc. | sufocar; (gás) pode causar asfixia e envenenamento}
  \end{Phonetics}
\end{Entry}

\begin{Entry}{熏香}{14,9}{⽕,⾹}
  \begin{Phonetics}{熏香}{xun1xiang1}
    \definition{s.}{incenso}
    \definition{v.}{exalar fumaça; fumigar}
  \end{Phonetics}
\end{Entry}

\begin{Entry}{熏陶}{14,10}{⽕,⾩}
  \begin{Phonetics}{熏陶}{xun1tao2}[][HSK 7-9]
    \definition{v.}{edificar; nutrir; exercer uma influência edificante sobre; as pessoas com quem você convive ao longo do tempo exercem gradualmente uma influência positiva em seu estilo de vida, pensamentos, comportamento, caráter e conhecimento}
  \synonymref{教会}{jiao1hui4}
  \synonymref{教导}{jiao4dao3}
  \synonymref{教授}{jiao4shou4}
  \synonymref{教学}{jiao4xue2}
  \synonymref{教养}{jiao4yang3}
  \synonymref{教育}{jiao4yu4}
  \synonymref{陶冶}{tao2ye3}
  \end{Phonetics}
\end{Entry}

%%%%%%%%%% 熙 %%%%%%%%%%
\subsection*{熙}\addcontentsline{loh}{figure}{熙}

\begin{Entry}{熙}{14}{⽕}
  \begin{Phonetics}{熙}{xi1}
    \definition{adj.}{Literário: brilhante; ensolarado | Literário: próspero |Literário: feliz; alegre; divertido}
    \definition{v.}{expor ao sol | levantar"-se;  ascender | jogar; brincar}
  \end{Phonetics}
\end{Entry}

\begin{Entry}{熙来攘往}{14,7,20,8}{⽕,⽊,⼿,⼻}
  \begin{Phonetics}{熙来攘往}{xi1lai2-rang3wang3}
    \definition{expr.}{a agitação de grandes multidões; o vai e vem em meio a multidões; repleto de atividade; pessoas circulando freneticamente}
  \synonymref{车水马龙}{che1shui3ma3long2}
  \synonymref{川流不息}{chuan1liu2-bu4xi1}
  \synonymref{络绎不绝}{luo4yi4-bu4jue2}
  \synonymref{熙熙攘攘}{xi1xi1-rang3rang3}
  \end{Phonetics}
\end{Entry}

\begin{Entry}{熙熙攘攘}{14,14,20,20}{⽕,⽕,⼿,⼿}
  \begin{Phonetics}{熙熙攘攘}{xi1xi1-rang3rang3}[][HSK 7-9]
    \definition{expr.}{``Repleto de atividade.''; pessoas circulando freneticamente}
  \seealsoref{熙来攘往}{xi1lai2-rang3wang3}
  \synonymref{车水马龙}{che1shui3ma3long2}
  \synonymref{熙来攘往}{xi1lai2-rang3wang3}
  \end{Phonetics}
\end{Entry}

%%%%%%%%%% 熬 %%%%%%%%%%
\subsection*{熬}\addcontentsline{loh}{figure}{熬}

\begin{Entry}{熬}{14}{⽕}
  \begin{Phonetics}{熬}{ao1}
    \definition{v.}{ensopar; ferver; cozinhar em água}
  \synonymref{挨}{ai2}
  \synonymref{炖}{dun4}
  \synonymref{忍}{ren3}
  \synonymref{煮}{zhu3}
  \antonymref{享受}{xiang3shou4}
  \end{Phonetics}
  \begin{Phonetics}{熬}{ao2}[][HSK 7-9]
    \definition{v.}{ferver; ensopar; fazer uma decocção; cozinhar em fogo baixo por muito tempo | preparar; infundir; extrair a essência por fervura longa | resistir; suportar (angústia, tempos difíceis, etc.)}
  \end{Phonetics}
\end{Entry}

\begin{Entry}{熬夜}{14,8}{⽕,⼣}
  \begin{Phonetics}{熬夜}{ao2/ye4}[][HSK 7-9]
    \definition{v.+compl.}{ficar acordado a noite toda ou até tarde da noite}
  \synonymref{通宵}{tong1xiao1}
  \end{Phonetics}
\end{Entry}

%%%%%%%%%% 熟 %%%%%%%%%%
\subsection*{熟}\addcontentsline{loh}{figure}{熟}

\begin{Entry}{熟}{15}{⽕}
  \begin{Phonetics}{熟}{shou2}
    \definition{adj.}{sinônimo de 熟 \dpy{shu2}, que significa cozido, maduro ou familiar}
  \antonymref{生}{sheng1}
  \end{Phonetics}
  \begin{Phonetics}{熟}{shu2}[][HSK 2]
    \definition{adj.}{maduro (frutos) | pronto; cozido | processado, fabricado ou exercitado | familiar, bem conhecido; conhecido por ser comum ou frequentemente utilizado | habilidoso;  (trabalho, tecnologia) experiente; não é novato | profundo; sólido}
  \antonymref{浅}{qian3}
  \antonymref{生}{sheng1}
  \end{Phonetics}
\end{Entry}

\begin{Entry}{熟人}{15,2}{⽕,⼈}
  \begin{Phonetics}{熟人}{shu2ren2}[][HSK 3]
    \definition[位,名,个,些]{s.}{amigo; conhecido; pessoas que se conhecem há muito tempo; pessoas que são muito familiares}
  \end{Phonetics}
\end{Entry}

\begin{Entry}{熟练}{15,8}{⽕,⽷}
  \begin{Phonetics}{熟练}{shu2lian4}[][HSK 4]
    \definition{adj.}{especializado; proficiente; qualificado; habilidoso}
  \end{Phonetics}
\end{Entry}

\begin{Entry}{熟悉}{15,11}{⽕,⼼}
  \begin{Phonetics}{熟悉}{shu2xi5}[][HSK 5]
    \definition{adj.}{familiarizado com; não ser estranho}
    \definition{v.}{estar familiarizado com; saber claramente que | conhecer bem algo ou alguém; compreender e dominar (a situação) através da observação ou da experiência}
  \end{Phonetics}
\end{Entry}

%%%%%%%%%% 燃 %%%%%%%%%%
\subsection*{燃}\addcontentsline{loh}{figure}{燃}

\begin{Entry}{燃}{16}{⽕}
  \begin{Phonetics}{燃}{ran2}
    \definition{v.}{queimar | acender; inflamar}
  \end{Phonetics}
\end{Entry}

\begin{Entry}{燃气}{16,4}{⽕,⽓}
  \begin{Phonetics}{燃气}{ran2qi4}[][HSK 7-9]
    \definition{s.}{combustível gasoso, como gás de carvão, biogás, gás natural e gás liquefeito de petróleo}
  \end{Phonetics}
\end{Entry}

\begin{Entry}{燃放}{16,8}{⽕,⽅}
  \begin{Phonetics}{燃放}{ran2fang4}[][HSK 7-9]
    \definition{v.}{acender (fogos de artifício, etc.); acender fogos de artifício, etc., para causar uma explosão}
  \end{Phonetics}
\end{Entry}

\begin{Entry}{燃油}{16,8}{⽕,⽔}
  \begin{Phonetics}{燃油}{ran2you2}[][HSK 7-9]
    \definition{s.}{óleo combustível}[燃油价格持续上涨了。===Os preços dos combustíveis continuaram a subir.]
  \end{Phonetics}
\end{Entry}

\begin{Entry}{燃料}{16,10}{⽕,⽃}
  \begin{Phonetics}{燃料}{ran2liao4}[][HSK 4]
    \definition[种]{s.}{combustível; carburante; substâncias que podem gerar calor e energia luminosa quando queimadas podem ser divididas em três tipos de acordo com sua forma: combustível sólido (como carvão, carvão vegetal, madeira), combustível líquido (como gasolina, querosene) e combustível gasoso (como gás de carvão, biogás); também se refere a substâncias que podem gerar energia nuclear, como urânio, plutônio, etc.}
  \end{Phonetics}
\end{Entry}

\begin{Entry}{燃烧}{16,10}{⽕,⽕}
  \begin{Phonetics}{燃烧}{ran2shao1}[][HSK 4]
    \definition{v.}{queimar; acender | arder; inflamar; ferver; metáfora para as emoções de uma pessoa serem muito fortes, como um fogo ardente}
  \end{Phonetics}
\end{Entry}

%%%%%%%%%% 燕 %%%%%%%%%%
\subsection*{燕}\addcontentsline{loh}{figure}{燕}

\begin{Entry}{燕}{16}{⽕}
  \begin{Phonetics}{燕}{yan1}
    \definition*{s.}{Hebei do Norte | Yan, um estado vassalo de Zhou nas atuais províncias de Hebei e Liaoning | Sobrenome: Yan}
  \end{Phonetics}
  \begin{Phonetics}{燕}{yan4}
    \definition{s.}{andorinha (uma família de pássaros) | conforto; desfruta | festa; banquete; jantar}
  \end{Phonetics}
\end{Entry}

\begin{Entry}{燕子}{16,3}{⽕,⼦}
  \begin{Phonetics}{燕子}{yan4zi5}[][HSK 7-9]
    \definition[只,窝]{s.}{andorinha}
  \end{Phonetics}
\end{Entry}

%%%%%%%%%% 爆 %%%%%%%%%%
\subsection*{爆}\addcontentsline{loh}{figure}{爆}

\begin{Entry}{爆}{19}{⽕}
  \begin{Phonetics}{爆}{bao4}[][HSK 6]
    \definition{v.}{explodir; estourar | fritar rapidamente; ferver rapidamente | aparecer (ou ocorrer) inesperadamente}
  \seealsoref{爆炸}{bao4zha4}
  \end{Phonetics}
\end{Entry}

\begin{Entry}{爆发}{19,5}{⽕,⼜}
  \begin{Phonetics}{爆发}{bao4fa1}[][HSK 6]
    \definition{v.}{entrar em erupção; explodir | estourar; irromper; ocorrer de forma repentina e violenta}
  \synonymref{发作}{fa1zuo4}
  \antonymref{发动}{fa1dong4}
  \end{Phonetics}
\end{Entry}

\begin{Entry}{爆竹}{19,6}{⽕,⽵}
  \begin{Phonetics}{爆竹}{bao4zhu2}[][HSK 7-9]
    \definition[串,个]{s.}{fogo de artifício}
  \end{Phonetics}
\end{Entry}

\begin{Entry}{爆米花}{19,6,7}{⽕,⽶,⾋}
  \begin{Phonetics}{爆米花}{bao4mi3hua1}
    \definition[袋]{s.}{pipoca (de milho); pipoca de arroz; alimentos feitos pelo aquecimento e expansão de arroz, grãos de milho, etc.}
  \end{Phonetics}
\end{Entry}

\begin{Entry}{爆冷门}{19,7,3}{⽕,⼎,⾨}
  \begin{Phonetics}{爆冷门}{bao4leng3men2}[][HSK 7-9]
    \definition{s.}{um avanço | uma reviravolta (especialmente nos esportes) | reviravolta inesperada dos acontecimentos}
    \definition{v.}{dar um golpe}
  \end{Phonetics}
\end{Entry}

\begin{Entry}{爆炸}{19,9}{⽕,⽕}
  \begin{Phonetics}{爆炸}{bao4zha4}[][HSK 6]
    \definition{s.}{explosão}
    \definition{v.}{explodir; explodir; detonar | aumentar bruscamente em um curto espaço de tempo (de quantidade)}
  \seealsoref{爆}{bao4}
  \end{Phonetics}
\end{Entry}

\begin{Entry}{爆满}{19,13}{⽕,⽔}
  \begin{Phonetics}{爆满}{bao4man3}[][HSK 7-9]
    \definition{v.}{(teatro, cinema, estádio, etc.) lotar; ter casa cheia | estar lotado}
  \antonymref{缺失}{que1shi1}
  \end{Phonetics}
\end{Entry}

%%%%% EOF %%%%%

